\begin{abstract}
Large-scale unsupervised language models (LMs) acquire extensive world knowledge and some inferential abilities through their training, yet exerting fine-grained control over their outputs remains challenging because their training paradigm is entirely unsupervised. Methods to address this limitation typically involve obtaining human annotations that express preferences between different model outputs, then adjusting the LM via fine-tuning to conform more closely to these preferences. This adjustment is often accomplished through reinforcement learning from human feedback (RLHF), a procedure that is both intricate and prone to instability; RLHF requires learning a reward model that encodes human preference judgments, then training the unsupervised LM through reinforcement learning to optimize this reward signal while maintaining similarity to the initial model. In this work, we propose an alternative reward model parameterization within RLHF that admits an analytical solution for the corresponding optimal policy, thereby reframing the conventional RLHF objective into a straightforward classification loss. Our approach, termed \textit{Direct Preference Optimization} (DPO), yields an algorithm that is stable, efficient, and computationally inexpensive, obviating the necessity for language model sampling during fine-tuning and reducing dependence on extensive hyperparameter tuning. Empirical results demonstrate that DPO matches or surpasses state-of-the-art alignment techniques in steering LMs according to human preferences. In particular, DPO outperforms PPO-based RLHF in modulating sentiment and achieves comparable or superior results in summarization and single-turn dialogue, all while offering significant simplicity in implementation and training.
\end{abstract}

\section{Introduction}
Large unsupervised language models (LMs) trained on extensive datasets have demonstrated unexpected abilities~\citep{chowdhery2022palm, brown2020language, touvron2023llama, bubeck2023sparks}. Nevertheless, the training data for these models is generated by humans whose goals, motivations, and expertise can vary widely. Inevitably, not all of these qualities are ideal to emulate; for instance, though an AI programming assistant should \textit{recognize} typical coding errors to address them effectively, we ultimately wish for the model to favor the (often uncommon) examples of high-quality programming contained in its training corpus during generation. In a similar vein, it is beneficial if a language model is \textit{familiar} with widely held misconceptions, but it is undesirable for the model to present such misconceptions as accurate in a corresponding fraction of its outputs. Thus, being able to distinguish and elicit the \emph{preferred outputs and behaviors} from the broader spectrum of its \textit{knowledge and skills} is essential for developing AI systems that are robust, high-performing, and controllable \citep{ouyang2022training}. Although reinforcement learning (RL) has been the dominant approach to aligning LMs with human preferences, we will demonstrate that the RL objective typically used in current frameworks can in fact be solved exactly using a straightforward binary cross-entropy loss, which streamlines the entire preference learning process.

Broadly, current strategies for imparting desired traits in language models depend on curated human preference datasets that define the types of outputs regarded as safe and beneficial. This preference alignment phase is conducted following the model’s initial large-scale unsupervised pre-training over massive text datasets. While supervised fine-tuning on human-generated, high-quality response samples provides a simple route to preference learning, the most effective methods are based on reinforcement learning from human (or AI) feedback (RLHF/RLAIF; \citep{christiano2017deep, bai2022constitutional}). These RLHF approaches involve training a reward model on human preference data and subsequently applying RL to adjust the language model’s output policy to prioritize responses with high reward scores, while ensuring the adapted policy remains sufficiently close to the original. Although RLHF yields language models with remarkable performance in conversation and coding tasks, its training workflow is notably more complicated than supervised approaches, requiring multiple LMs and in-the-loop policy sampling, thus demanding substantial computational resources.

In this work, we present a method for aligning language models with human preferences without resorting to explicit reward models or reinforcement learning frameworks. We introduce \textit{{Direct Preference Optimization} (DPO)}, an approach that implicitly targets the same reward maximization objective with a KL-divergence regularization as conventional RLHF techniques, but which offers a more straightforward implementation and training procedure. Conceptually, DPO adjusts the relative log likelihood of favored versus unfavored responses, incorporating an adaptive, instance-specific weighting mechanism that mitigates the risk of model collapse observed when using a basic probability ratio objective. As with prior work, DPO leverages a foundational preference model (such as the Bradley-Terry model; \cite{bradley1952rankanalysis}) to evaluate the degree of agreement between a reward signal and observed preference labels. In contrast to existing strategies—which define a preference loss to optimize a reward model before training a policy against that reward—DPO directly re-parameterizes the preference loss in terms of the policy itself. This enables DPO to optimize policy parameters with a straightforward binary cross-entropy objective using datasets of pairwise human preferences, thereby yielding an optimal policy for a reward function implicitly learned from the preference data.

The primary contribution of our study is the introduction of Direct Preference Optimization (DPO), an uncomplicated and reinforcement learning-free method for preference-driven language model training. Through empirical evaluations, we demonstrate that DPO matches or surpasses the performance of established approaches, such as PPO-based RLHF, in various domains like sentiment control, text summarization, and conversational modeling, and it scales effectively to language models with up to 6B parameters.

\section{Related Work}

Self-supervised language models that scale to larger sizes exhibit the capability to address certain tasks in a zero-shot manner \citep{radford2019language} or by leveraging few-shot prompt examples \citep{gpt3,megatron,chowdhery2022palm}. Nevertheless, augmenting these models’ downstream performance and improving alignment with user intent can be achieved through fine-tuning on curated datasets composed of instructions and human-provided completions \citep{mishra-etal-2022-cross,sanh2022multitask,chung2022scaling,thoppilan2022lamda}. The process of `instruction-tuning' facilitates LLMs in adapting to novel instructions beyond those presented during fine-tuning, enhancing their general applicability and usefulness \citep{chung2022scaling}. Although instruction tuning has proven effective, it is generally easier to obtain \textit{relative} assessments of response quality from non-expert humans than to acquire expert-annotated demonstrations. Consequently, recent research has pursued fine-tuning LLMs using datasets structured around human preferences, yielding advancements in translation \citep{kreutzer-etal-2018-reliability}, summarization \citep{stiennon2022learning,ziegler2020finetuning}, story generation \citep{ziegler2020finetuning}, and instruction-following abilities \citep{ouyang2022training,ramamurthy2023is}. This line of work typically entails first fitting a neural network-based reward model to align with human preference data under a preference framework such as the Bradley-Terry model \citep{bradley1952rankanalysis}, then adapting the language model by maximizing this learned reward via reinforcement learning approaches, with popular choices including REINFORCE \citep{williams1992reinforce}, proximal policy optimization (PPO; \cite{schulman2017proximal}), and related variants \citep{ramamurthy2023is}. Parallel efforts involve the use of instruction-following LLMs—fine-tuned using human feedback—to synthesize further preference data that targets attributes like safety or harmlessness \citep{bai2022constitutional}, relying on minimal direct supervision and instead employing a textual rubric to guide the LLM’s annotation process. These developments epitomize the intersection between two strands of research: one focused on reinforcement learning-based language model training for diverse objectives~\citep{Ranzato2015SequenceLT,paulus2018a,wu2018learning}, and the other on generalizable strategies for incorporating human preferences into model learning \citep{christiano2017deep,kupcsik2018learning}. Yet, despite the attractiveness of harnessing relative human preferences, implementing reinforcement learning-based fine-tuning for large-scale language models poses substantial practical difficulties; the present work proposes a theoretically-grounded methodology for preference optimization that circumvents the need for reinforcement learning.

Beyond applications in language, the problem of learning policies from preferences has been explored within both bandit and reinforcement learning frameworks, resulting in a variety of proposed solutions. In the contextual bandit scenario, utilizing preferences or rankings rather than explicit rewards leads to the framework of contextual dueling bandits (CDB; \cite{yue2012karmed,dudik2015contextual}). When absolute rewards are unavailable, the theoretical framework for CDBs replaces the standard notion of an optimal policy with a \textit{von Neumann winner}—a policy that is expected to achieve at least a 50\% win rate in head-to-head comparisons against \textit{any} alternative policy \citep{dudik2015contextual}. A notable distinction, however, is that CDBs operate in an online manner, continually receiving new preference feedback, whereas human preference-based learning typically utilizes a static, offline collection of annotated action pairs \citep{yan2022human}. In a related vein, \textit{preference-based RL} (PbRL) focuses on learning from binary preferences inferred via an \textit{unknown} 'scoring' function instead of observed rewards \citep{BusaFekete2014,ruiz2023dueling}. Numerous approaches for PbRL have been introduced, including those designed to leverage off-policy preference information; these typically proceed by first constructing an explicit model of the latent scoring (reward) function before optimizing the policy accordingly \citep{jain2013learning,BusaFekete2014,christiano2017deep,sadigh2017active,kupcsik2018learning}. Departing from this two-stage process, we introduce a unified approach that learns a policy directly optimized to align with preference data.

\section{Preliminaries}\label{section:prelims}

Here, we summarize the RLHF pipeline as outlined in \citeauthor{ziegler2020finetuning} (and expanded in subsequent works such as \citep{stiennon2022learning, bai2022training, ouyang2022training}). This procedure is generally divided into three sequential stages: (1) supervised fine-tuning (SFT); (2) reward modelling through preference elicitation and annotation; and (3) reinforcement learning-based optimization.

\textbf{SFT}: The RLHF process initiates by applying supervised learning to adapt a pre-trained language model using curated, task-specific datasets (for dialogue, summarization, etc.), resulting in an initial model $\pi^\text{SFT}$. 

\textbf{Reward Modelling Phase}: Next, prompts $x$ are fed into the SFT model to generate paired candidate outputs $(y_1, y_2)\sim \pi^\text{SFT}(y \mid x)$. These output pairs are evaluated by human annotators, who select the preferred response, formally $y_w\succ y_l \mid x$, where $y_w$ and $y_l$ indicate the responses judged preferable and less preferable, respectively. The observed preferences are presumed to originate from an unobserved, latent reward function $r^*(y, x)$. Several probabilistic models exist to capture these preferences, with the Bradley-Terry (BT) model \cite{bradley1952rankanalysis} being widely utilized; alternatively, the more flexible Plackett-Luce models \citep{plackett1975analysis, luce2012individual} are suitable when multiple candidate rankings are available. According to the BT model, the human preference probability distribution $p^*$ can be specified as follows:

\begin{equation}\label{eq:bradley-terry}
    p^*(y_1\succ y_2 \mid x)=\frac{\exp\left(r^*(x, y_1)\right)}{\exp\left(r^*(x, y_1)\right) + \exp\left(r^*(x, y_2)\right)}.
\end{equation}

Suppose we are given a fixed set of comparison data, denoted as $\mathcal{D} = \bigl\{x^{(i)}, y_w^{(i)}, y_l^{(i)}\bigr\}_{i=1}^N$, drawn from the underlying distribution $p^*$. We can then introduce a parametric reward function $r_{\phi}(x, y)$ and fit its parameters using the maximum likelihood principle. When cast as a binary classification task, the associated loss is the negative log-likelihood:

\begin{equation}\label{eq:reward_model}
    \mathcal{L}_R(r_{\phi}, \mathcal{D}) = -\mathbb{E}_{(x, y_w, y_l)\sim \mathcal{D}}\bigl[\log \sigma(r_{\phi}(x, y_w)- r_{\phi}(x, y_l))\bigr]
\end{equation}

where $\sigma$ denotes the logistic sigmoid. In large language model (LM) applications, $r_{\phi}(x, y)$ is frequently initialized with parameters from the SFT model $\pi^\text{SFT}(y \mid x)$, augmented by a linear projection atop the final transformer layer to output a scalar reward \cite{ziegler2020finetuning}. To minimize variance in the learned reward function, normalization is typically performed such that  $\mathbb{E}_{x,y\sim \mathcal{D}}\left[r_\phi(x, y)\right] = 0$ for all $x$.

\textbf{RL Fine-Tuning Phase}: During the reinforcement learning stage, this reward function supplies supervisory signals to guide the language model. Building upon prior works~\citep{jaques2017sequence, jaques2020human}, policy optimization is cast as

\begin{equation}\label{eq:RL}
\max_{\pi_{\theta}}  \mathbb{E}_{x\sim \mathcal{D}, y\sim \pi_{\theta}(y \mid x)}\bigl[r_{\phi}(x, y)\bigr] - \beta\mathbb{D}_{\textrm{KL}}\bigl[\pi_{\theta}(y\mid x)\mid \mid \pi_\text{ref}(y\mid x)\bigr],
\end{equation}

where the hyperparameter $\beta$ controls how much the learned policy can diverge from the base (reference) policy $\pi_\text{ref}$, typically set to the SFT model $\pi^\text{SFT}$. The policy $\pi_\theta$ is also commonly initialized from $\pi^\text{SFT}$. Including this regularization term serves two key purposes: constraining policy updates to regions where the reward model remains valid, and preserving output diversity to avoid the collapse onto narrow, high-reward outputs. Because language generation is inherently discrete, the resulting objective cannot be differentiated directly, and reinforcement learning methods are required to optimize it. A standard practice \citep{ziegler2020finetuning, stiennon2022learning, bai2022training, ouyang2022training} is to define the reward as ${r(x, y) = r_{\phi}(x, y) -\beta (\log \pi_{\theta}(y\mid x) - \log \pi_\text{ref}(y\mid x))}$, then employ the PPO algorithm \cite{schulman2017proximal} for maximization.

\section{Direct Preference Optimization}\label{sec:DPO}

Given the computational complexity of applying reinforcement learning to large-scale models such as LMs, we seek a streamlined alternative for policy learning that works directly with preference data. Departing from standard RLHF paradigms—which first fit an explicit reward model and subsequently optimize it via RL—our strategy utilizes a specific reward model structure that admits a closed-form solution for its optimal policy, eliminating the necessity for a reinforcement learning training loop. As we elaborate in the following sections, the core idea is to exploit an analytic correspondence between reward functions and their induced optimal policies. This enables reparametrization: reformulating a loss defined on the reward space as an equivalent loss on the space of policies. In this framework, there is no separate reward model; instead, the policy network itself encapsulates both the language model and its (

\textbf{Deriving the DPO objective.} We begin with the RL objective described in earlier studies, Eq.~\ref{eq:RL}, under a general reward formulation $r$. As established by previous work~\citep{peters2007reinforcement, peng2019advantage, korbak2022reinforcement, go2023aligning}, it can be directly shown that the solution to the KL-regularized reward maximization in Eq.~\ref{eq:RL} is:

\begin{equation}\label{eq:op_policy}
    \pi_r(y\mid x) = \frac{1}{Z(x)}\pi_\text{ref}(y\mid x)\exp\left(\frac{1}{\beta}r(x, y)\right),
\end{equation}

where the normalization constant $Z(x) =\sum_{y}\pi_\text{ref}(y\mid x)\exp\left(\frac{1}{\beta}r(x, y)\right)$. The full derivation is available in Appendix \ref{app:derivation1}. Notably, even with the maximum likelihood estimate $r_{\phi}$ of the actual reward function $r^*$, direct computation of $Z(x)$ remains computationally challenging \citep{korbak2022reinforcement, go2023aligning}, which complicates practical use. Nevertheless, Eq.~\ref{eq:op_policy} can be manipulated to describe the reward as a function of the optimal policy $\pi_r$, the reference policy $\pi_\text{ref}$, and $Z(\cdot)$. Taking the logarithm of both sides of Eq.~\ref{eq:op_policy} and rearranging yields:

\begin{equation}\label{eq:main_eq}
    r(x,y) =\beta \log \frac{\pi_r(y\mid x)}{\pi_\text{ref}(y\mid x)} + \beta \log Z(x).
\end{equation}

This reformulation applies equally well to the true reward $r^*$ and its associated optimal policy $\pi^*$. Importantly, the Bradley-Terry framework relies exclusively on the difference between the rewards of two possible outputs; that is, ${p^*(y_1 \succ y_2 \mid x) = \sigma(r^*(x, y_1) - r^*(x, y_2))}$. By inserting the form from Eq.~\ref{eq:main_eq} for $r^*(x,y)$ into the preference model of Eq.~\ref{eq:bradley-terry}, the partition term disappears, enabling the human preference probability to be written purely in terms of $\pi^*$ and $\pi_\text{ref}$. Thus, the optimal RLHF policy $\pi^*$ following the Bradley-Terry setup meets the preference criteria:

\begin{equation}\label{eq:objective}
    p^*(y_1\succ y_2 \mid x)=\frac{1}{1 + \exp\left(\beta \log \frac{\pi^*(y_2\mid x)}{\pi_\text{ref}(y_2\mid x)} - \beta \log \frac{\pi^*(y_1\mid x)}{\pi_\text{ref}(y_1\mid x)}\right)}
\end{equation}

A step-by-step derivation appears in Appendix~\ref{app:derivation2}. Though Eq.~\ref{eq:objective} leverages the Bradley-Terry approach, analogous derivations apply for broader preference models such as Plackett-Luce~\citep{plackett1975analysis, luce2012individual}, details of which are provided in Appendix~\ref{app:plackett_luce_models}.

With this formulation, the probability of observed human preferences depends directly on the optimal policy, bypassing explicit reward modeling. As a result, one can define a maximum likelihood learning criterion for a parametrized policy $\pi_\theta$. Mirroring the approach used for reward modeling (cf. Eq.~\ref{eq:reward_model}), the policy learning objective becomes:

\begin{equation}\label{eq:optimum_model}
    \mathcal{L}_\text{DPO}(\pi_{\theta}; \pi_\text{ref}) = -\mathbb{E}_{(x, y_w, y_l)\sim \mathcal{D}}\left[\log \sigma \left(\beta \log \frac{\pi_{\theta}(y_w\mid x)}{\pi_\text{ref}(y_w\mid x)} - \beta \log \frac{\pi_{\theta}(y_l\mid x)}{\pi_\text{ref

Through this approach, we approximate an implicit reward via an alternative parameterization, where the corresponding optimal policy aligns with $\pi_\theta$. Furthermore, as our method can be viewed as fitting a reparametrized Bradley-Terry model, it inherits certain theoretical guarantees, such as consistency under appropriate assumptions on the distribution of preference data \cite{bong2022generalized}. We further elaborate on the theoretical underpinnings of DPO and its connections to related frameworks in Section~\ref{sec:theory}.

\textbf{How does the DPO update function?} To gain insight into the inner workings of DPO, one may analyze the gradient of the loss function $\mathcal{L}_\text{DPO}$. The gradient with respect to model parameters $\theta$ takes the following form:

\begin{multline*}\label{eq:gradient}
    \nabla_\theta \mathcal{L}_\text{DPO}(\pi_\theta;\pi_\text{ref}) = \\ -\beta\mathbb{E}_{(x, y_w, y_l) \sim \mathcal{D}} \bigg[\underbrace{\sigma(\hat{r}_\theta(x, y_l) - \hat{r}_\theta (x, y_w))}_\text{higher weight when reward estimate is wrong}\bigg[\underbrace{\nabla_\theta\log \pi(y_w \mid x)}_\text{increase likelihood of $y_w$} - \underbrace{\nabla_\theta\log\pi(y_l \mid x)}_\text{decrease likelihood of $y_l$}\bigg]\bigg],
\end{multline*}

Here, $\hat{r}_\theta(x, y) = \beta \log \frac{\pi_\theta(y \mid x)}{\pi_\text{ref}(y \mid x)}$ denotes the reward that is implicitly characterized by both the learned language model $\pi_\theta$ and the reference policy $\pi_\text{ref}$ (see Section~\ref{sec:theory} for further details). Conceptually, this gradient updates the model to increase the probability assigned to preferred completions $y_w$ while suppressing the probability of dispreferred completions $y_l$. Crucially, the contribution of each example is weighted according to the extent to which the implicit reward model $\hat{r}_\theta$ overestimates the less-preferred completions relative to the preferred ones—scaled by $\beta$. In other words, the more the reward model erroneously ranks $y_l$ above $y_w$, the greater its influence, reflecting the severity of the KL regularization. Our empirical results underscore the significance of this weighting scheme: omitting the weighting coefficient, as in a naive implementation, can lead to model collapse, as demonstrated in Appendix Table~\ref{tab:unlikelihood_generations}.

\textbf{DPO overview.} The standard DPO process proceeds as follows: First, for each prompt $x$, generate completion samples $y_1, y_2 \sim \pi_\text{ref}(\cdot \mid x)$, then assign human preference labels to assemble an offline preference dataset $\mathcal{D} = \{x^{(i)}, y_w^{(i)}, y_l)^{(i)}\}_{i=1}^N$. Second, update the language model $\pi_\theta$ by minimizing the objective $\mathcal{L}_\text{DPO}$, conditioned on $\pi_\text{ref}$, the dataset $\mathcal{D}$, and a chosen $\beta$ value.  
In real-world scenarios, it is preferable to utilize existing public preference datasets rather than collect fresh human annotations and new samples. These datasets typically originate from $\pi^\text{SFT}$-sampled completions, prompting the initialization $\pi_\text{ref} = \pi^\text{SFT}$ where feasible. If $\pi^\text{SFT}$ is unavailable, $\pi_\text{ref}$ is initialized via maximizing the likelihood of the preferred completions ${(x, y_w)}$ over the dataset, i.e., setting ${\pi_\text{ref} = \argmax_{\pi}\mathbb{E}_{x, y_w \sim \mathcal{D}}\left[\log \pi(y_w \mid x)\right]}$. This approach helps address the distribution gap between the ideal but unattainable reference distribution and the actual $\pi_\text{ref}$ employed in DPO. More granular implementation specifics and information on hyperparameter settings are provided in Appendix~\ref{app:implementation}.

\section{Theoretical Analysis of DPO}
Here, we develop a deeper theoretical understanding of the DPO algorithm, offering formal justification and clarifying how DPO avoids the pitfalls found in actor-critic frameworks commonly applied in RLHF (such as PPO~\cite{schulman2017proximal}).

\label{sec:theory}

\subsection{Your Language Model Is Secretly a Reward Model} DPO achieves both implicit reward learning and policy optimization through a unified maximum likelihood approach, removing the need for explicit reward modeling and reinforcement learning procedures. The objective function (Eq. \ref{eq:main_eq}) is mathematically equivalent to a Bradley-Terry framework, with rewards given by $r^*(x, y) = \beta \log\frac{\pi^*_\theta(y \mid x)}{\pi_\text{ref}(y \mid x)}$. Here, we directly update the parametric policy $\pi_{\theta}$, mirroring the reward model training specified in Eq. \ref{eq:reward_model} but expressed under a change of variables. This section systematically constructs the rationale for this transformation, demonstrates that it does not impose limitations on the spectrum of learnable reward functions, and shows how it enables exact retrieval of the optimal policy. We begin by introducing a formal equivalence between reward functions.

\begin{definition}
Two reward functions $r(x, y)$ and $r'(x, y)$ are said to be equivalent if and only if ${r(x, y)-r'(x, y) = f(x)}$ for some function $f$.
\end{definition}

This notion of equivalence constitutes an equivalence relation, thereby dividing the space of reward functions into distinct equivalence classes. We summarize two immediate consequences below:

\begin{lemma}\label{lemma:same_prefrence}
Within the Plackett-Luce framework, and particularly for the Bradley-Terry model, any pair of reward functions belonging to the same equivalence class will induce identical preference distributions.
\end{lemma}

\begin{lemma}\label{lemma:same_policy}
    If two reward functions are members of the same equivalence class, they will result in the same optimal policy in the context of the constrained RL problem.
\end{lemma}

Since the proofs follow straightforwardly, they are postponed to Appendix \ref{app:lemma1}. The first lemma reflects a widely-recognized limitation of the Plackett-Luce family—specifically, the under-specification of these models \cite{plackett1975analysis}. As a consequence, it is customary to impose additional identifiability constraints in order to obtain meaningful MLE solutions from Eq.~\ref{eq:reward_model} \cite{bong2022generalized}. The second lemma highlights that all reward functions within an equivalence class prescribe the same optimal policy, which means our ultimate goal reduces to identifying any member of the optimal class of reward functions. The theorem below, which we prove in Appendix~\ref{app:thm1}, formalizes this characterization:

\begin{theorem}\label{thm:main}
    Assuming mild conditions, every reward equivalence class compatible with the Plackett-Luce (or Bradley-Terry) framework admits a representation of the form ${r(x, y) = \beta \log \frac{\pi(y\mid x)}{\pi_\text{ref}(y\mid x)}}$, where $\pi(y\mid x)$ denotes some model, and $\pi_\text{ref}(y \mid x)$ is a fixed reference model.
\end{theorem}

\begin{sproof}
    Let $r(x, y)$ be an arbitrary reward function, with the associated optimal model $\pi_r(y \mid x)$ defined as in Eq.~\ref{eq:op_policy}. Our goal is to demonstrate that there exists a member in the equivalence class of $r$ expressible through the above reparameterization. Consider the mapping
\begin{equation}
    f(r; \pi_\text{ref}, \beta)(x, y) = r(x, y) - \beta\log\sum_{y}\pi_\text{ref}(y\mid x)\exp\left(\frac{1}{\beta}r(x, y)\right)
\end{equation}
This mapping $f$ adjusts $r(x, y)$ by subtracting the log-partition function, which depends only on $x$, thereby ensuring $f(r; \pi_\text{ref}, \beta)(x, y)$ resides in the same equivalence class as $r(x, y)$. By substituting $r$ with the expression on the right-hand side of Eq.~\ref{eq:main_eq}, which holds for arbitrary reward functions, it follows that $f(r; \pi_\text{ref}, \beta)(x, y) = \beta \log \frac{\pi_r(y\mid x)}{\pi_\text{ref}(y\mid x)}$. Hence, the mapping $f$ yields a representative of the equivalence class in the desired canonical form, validating the generality of our proposed reward model parameterization.
\end{sproof}

An alternative interpretation of Theorem~\ref{thm:main} is that it precisely characterizes the reward function chosen by the DPO reparameterization within each equivalence class, specifically the function satisfying:

\begin{equation}\label{eq:lag_p}
     \sum_{y}\underbrace{\pi_\text{ref}(y\mid x)\exp\left(\frac{1}{\beta}r(x, y)\right)}_{=\pi(y\mid x)\text{, using Thm.~\ref{thm:main} reparam.}} = 1,
\end{equation}

In other words, this condition ensures that $\pi(y\mid x)$ defines a legitimate probability distribution with non-negative values that sum to one. Furthermore, referencing Eq.~\ref{eq:op_policy}, we observe that Eq.~\ref{eq:lag_p} functions as the normalization constant—or partition function—for the optimal policy generated by the reward $r(x, y)$. The principal contribution of the DPO method lies in its ability to introduce specific constraints to the typically under-determined Plackett-Luce (and, by extension, Bradley-Terry) preference model class. This procedure retains the expressive capacity of the reward modeling class while simultaneously allowing the optimal policy of Eq.~\ref{eq:op_policy} to be computed explicitly for any prompt $x$.

\subsection{Instability of Actor-Critic Algorithms} 

The framework we develop can additionally shed light on the sources of instability often encountered with classical actor-critic approaches in RLHF, such as PPO. To do so, we focus on the RL fine-tuning stage delineated in Section \ref{section:prelims}, adopting the typical RLHF training protocol. Within this context, connections arise to the control-as-inference paradigm \cite{levine2018reinforcement} for constrained reinforcement learning problems, as formalized in \ref{eq:RL}. Assume a policy model parameterized as $\pi_{\theta}(y\mid x)$. The learning objective involves minimizing $\mathbb{D}_{\text{KL}}[\pi_{\theta}(y|x) \mid \mid \pi^*(y\mid x)]$, where $\pi^*$ denotes the optimal policy per Eq.~\ref{eq:optimum_model}, derived from the reward function $r_{\phi}(y, x)$. Through further manipulation, the following optimization target emerges:

\begin{equation}\label{eq:AC}
    \max_{\pi_{\theta}}\mathbb{E}_{\pi_{\theta}(y\mid x)}\bigg[\underbrace{r_{\phi}(x, y) -\beta\log\sum_{y}\pi_\text{ref}(y\mid x)\exp\left(\frac{1}{\beta}r_{\phi}(x, y)\right)}_{f(r_{\phi}, \pi_\text{ref}, \beta)} - \underbrace{\beta\log\frac{\pi_{\theta}(y\mid x)}{\pi_\text{ref}(y\mid x)}}_{\text{KL}}\bigg]
\end{equation}

The objective here aligns with that targeted in previous studies 
\citep{ziegler2020finetuning, stiennon2022learning, bai2022training, ouyang2022training}, where the DPO-equivalent reward is utilized for the reward class $r_{\phi}$. In this context, the normalization component within $f(r_{\phi}, \pi_\text{ref}, \beta)$ can be seen as representing the soft value function corresponding to the reference policy $\pi_\text{ref}$. Although this normalization factor does not alter the set of optimal policies, its omission could result in a policy gradient with elevated variance, thereby impeding the stability of the learning process. To handle this, one could introduce a learned value function to approximate the normalization, though this approach is often challenging to optimize. As an alternative, earlier methods have employed normalization via a human completion baseline, effectively yielding a one-sample Monte Carlo estimate of the normalizer. Unlike these approaches, the DPO reparameterization constructs a reward function that does not depend on any explicit baselines.

\section{Experiments}
This section presents empirical analyses of how effectively DPO can train policies from preference data. We begin by examining a controlled text-generation environment to explore the efficiency with which DPO balances reward maximization and KL-divergence minimization relative to reference policies, comparing its performance to established preference learning methods such as PPO. Subsequently, we assess DPO in the context of larger-scale models and more challenging RLHF tasks, encompassing both summarization and conversational dialogue. Our findings indicate that, even with minimal hyperparameter tuning, DPO matches or surpasses strong baselines like RLHF using PPO and outperforms approaches that select the best of $N$ samples as judged by a learned reward model. Preceding the main results, we detail the experimental procedures used; supplementary methodological specifics are available in Appendix~\ref{app:exp_details}.

\textbf{Tasks.} In our empirical studies, we investigate three open-ended text generation scenarios. Across all tasks, algorithms are trained to learn a policy from a dataset of human preference pairs, represented as $\mathcal{D}=\bigl\{x^{(i)}, y_w^{(i)}, y_l^{(i)}\bigr\}_{i=1}^N$. For \textbf{controlled sentiment generation}, the input $x$ consists of the beginning segment of a movie review selected from the IMDb dataset \cite{maas-EtAl:2011:ACL-HLT2011}, and the objective for the policy is to generate a continuation $y$ that expresses positive sentiment. To ensure objective evaluation, we synthesize preference pairs between candidate generations by employing a pre-trained sentiment classifier, assigning preferences such that $p(\text{positive}\mid x,y_w)>p(\text{positive}\mid x,y_l)$. For SFT, GPT-2-large is fine-tuned on the training subset of IMDb reviews until model performance converges (see further discussion in App~\ref{app:sentiment_details}). In the case of \textbf{summarization}, $x$ represents a Reddit forum post, with the policy tasked with producing a summary $y$ that captures the primary points. Following standard practice, we utilize the Reddit TL;DR summarization dataset \citep{volske-etal-2017-tl} supplemented with human preference data from \citeauthor{stiennon2022learning}. Here, an SFT model—fine-tuned on summaries authored by humans\footnote{\url{https://huggingface.co/CarperAI/openai_summarize_tldr_sft}}—is used in conjunction with RLHF via the TRLX framework \citep{leandro_von_werra_2023_7790115}. Notably, the human preference annotations were collected on outputs from a model closely related to, but not identical with, the SFT model deployed in this study. Lastly, for \textbf{single-turn dialogue}, each $x$ is a user prompt, ranging widely from science questions to personal advice inquiries. The policy must then provide a relevant and constructive response $y$. For this purpose, we draw on the Anthropic Helpful and Harmless dialogue dataset \citep{bai2022training}, which includes 170,000 dialogue instances between a user and an AI assistant, each concluding with two possible AI-generated replies and an accompanying human preference indicator for the superior response. In this experimental condition, a pre-trained SFT model is unavailable; thus, we create an SFT baseline by fine-tuning a standard language model solely on the preferred completions.

\textbf{Evaluation.} We employ two complementary evaluation strategies in our experiments. To investigate how effectively each algorithm can optimize the constrained reward maximization objective, we measure performance in the controlled sentiment generation scenario by plotting the frontier between achieved reward and KL-divergence from a reference policy. This is feasible because the underlying reward function (specifically, a sentiment classifier) is accessible in this setting. Conversely, in practical applications where the true reward function remains unknown, we turn to \textit{win rate} as an evaluation criterion, benchmarking each algorithm against a baseline policy. Here, GPT-4 serves as an automated surrogate for human assessment, evaluating both summary quality in summarization and helpfulness in single-turn dialogue. For summarization tasks, we designate the reference summary within the test set as the baseline; for dialogue, the baseline consists of the dataset’s preferred response. Although prior research indicates that LMs can outperform traditional metrics as evaluators \citep{Chen2023ExploringTU}, we supplement this with a human study to support our use of GPT-4 for scoring, as detailed in Sec.~\ref{sec:human-judgments}. Our results show a strong concordance between GPT-4 and human judgments, often matching or exceeding the level of agreement seen between different human annotators.

\textbf{Methods.} Alongside DPO, we assess several established methods for training language models in alignment with human preferences. For the summarization domain, we apply zero-shot prompting to \textbf{GPT-J} \citep{gpt-j}, and for dialogue, we utilize 2-shot prompting with \textbf{Pythia-2.8B} \citep{biderman2023pythia}. Our experiments also include the \textbf{SFT} approach and the \textbf{Preferred-FT} model, the latter of which is produced by fine-tuning via supervised learning on the selected completion $y_w$—using outputs from the SFT model in sentiment and summarization tasks, or from a generic language model for single-turn dialogue. Another supervised-like baseline is the \textbf{Unlikelihood}~\citep{welleck2019neural} objective, which encourages the model to boost the likelihood of $y_w$ while explicitly reducing the probability of $y_l$; we introduce an optional weighting factor $\alpha\in[0,1]$ for the 'unlikelihood' component. Additionally, we examine \textbf{PPO} \citep{schulman2017proximal}, using a reward function trained on preference annotations, as well as an oracle variant, \textbf{PPO-GT}, which leverages access to the actual reward function in the controlled sentiment context. In the sentiment generation experiments, two PPO-GT versions are tested: one is a standard off-the-shelf implementation \cite{leandro_von_werra_2023_7790115}, while the other incorporates reward normalization and fine-tuned hyperparameters for better outcomes; these enhancements are also applied when running standard PPO with learned rewards. Lastly, we include the \textbf{Best of $N$} strategy, which samples $N$ completions from the SFT model (or Preferred-FT for dialogue) and returns the response with the highest reward model score. Although this approach disentangles reward model accuracy from PPO optimization and demonstrates strong performance, it is computationally expensive—requiring $N$ samples for each test-time query—even for moderate values of $N$.

\subsection{How well can DPO optimize the RLHF objective?}

The reward maximization objective with a KL constraint, common in standard RLHF methods, is designed to strike a balance between optimizing for higher rewards and ensuring the learned policy does not diverge excessively from the reference policy. As such, when evaluating different methods, it is important to jointly consider the level of reward attained as well as the KL divergence from the reference; prioritizing higher rewards at the expense of a significant increase in KL is often suboptimal. In Figure~\ref{fig:frontier-tldr-main}, we present the trade-off frontier between reward and KL for several methods within the sentiment classification scenario. For each algorithm, multiple independent training experiments are conducted, each using a distinct policy conservativeness parameter (target KL $\in\{3,6,9,12\}$ for PPO, $\beta \in \{0.05,0.1,1,5\}$, $\alpha\in\{0.05,0.1,0.5,1\}$ for unlikelihood, different random seeds for preferred-FT), resulting in 22 total experimental runs. At intervals of 100 training steps up to convergence, we evaluate every policy on a fixed set of test prompts, reporting both the mean reward according to the actual reward function and the mean sequence-level KL\footnote{Specifically, this refers to summing the KL-divergences at each time step across the sequence.} to the reference policy, $\text{KL}\left(\pi\mid \mid \pi_\text{ref}\right)$. Our experiments demonstrate that DPO produces a reward-KL frontier that is considerably more efficient than the alternatives, attaining the highest rewards for a given KL, or lower KL for equivalent reward. This performance is striking for several reasons: firstly, even though both DPO and PPO seek to optimize the same formal objective, DPO displays a superior efficiency, clearly outperforming PPO along the reward/KL frontier. Secondly, DPO’s frontier is better than that of PPO even in the case where PPO has access to the exact, ground truth reward signals (PPO-GT).

\subsection{Can DPO scale to real preference datasets?}
\label{sec:dpo-real-datasets}
We proceed to assess the fine-tuning capabilities of DPO within the domains of summarization and single-turn dialogue tasks. In the context of summarization, commonly used automatic metrics like ROUGE have been shown to correlate weakly with human preferences~\citep{stiennon2022learning}, and previous studies indicate that leveraging PPO for fine-tuning language models based on human judgments can yield higher-quality summaries. To benchmark various techniques, we generate completions using the test partition of the TL;DR summarization dataset, measuring the average win rate relative to the reference completions within this test set. For all approaches, completions are sampled using temperature values between 0.0 and 1.0, and the corresponding win rates are presented in Figure~\ref{fig:frontier-tldr-main} (right). The fine-tuning for DPO, PPO, and Preferred-FT is conducted on the identical GPT-J SFT checkpoint\footnote{\url{https://huggingface.co/CarperAI/openai_summarize_tldr_sft}}. At a temperature setting of 0.0, DPO attains an approximate win rate of 61\%, surpassing PPO, which achieves about 57\% at its own optimal sampling temperature of 0.0. DPO also records a higher peak win rate compared with the best-of-$N$ baseline. It should be noted that we did not undertake extensive tuning of the $\beta$ hyperparameter for DPO, so these figures may not fully reflect its capacity. Additionally, DPO demonstrates substantially greater resilience to temperature variations than PPO, whose performance can decline to that of the untuned GPT-J model at elevated temperatures. Preferred-FT does not yield a notable performance gain relative to the SFT baseline. Further, a direct human evaluation comparison between DPO and PPO, detailed in Section~\ref{sec:human-judgments}, shows that DPO samples generated at temperature 0.25 were favored 58\% of the time over those produced by PPO at temperature 0.

For single-turn dialogue evaluation, we focus on the portion of the Anthropic HH dataset \citep{bai2022training} test set involving a single round of human-assistant interaction. In the GPT-4-based assessments, we utilize the test-preferred completions as references to calculate win rates across the different methods. Since there is no widely accepted SFT baseline for this scenario, our approach initiates with a pre-trained Pythia-2.8B model; we then apply Preferred-FT to construct a reference model using the selected completions to ensure outputs are representative of the target distribution, subsequently proceeding to train with DPO. 

To contextualize DPO's performance, we benchmark against two additional baselines: the optimal result among 128 Preferred-FT completions (noting that the Best of $N$ baseline performance saturates at $N=128$ for this task; refer to Appendix Figure~\ref{fig:best-of-n}), and a 2-shot prompt applied to the original Pythia-2.8B model. Our experiments reveal that DPO achieves performance matching or surpassing the best-performing temperature settings among these methods.

Additionally, we assess an RLHF model—trained using PPO on the Anthropic HH dataset\footnote{\url{https://huggingface.co/reciprocate/ppo_hh_pythia-6B}} and sourced from a reputable implementation\footnote{\url{https://github.com/CarperAI/trlx/tree/main/examples/hh}}—but find no combination of prompt or temperature yielding improvements over the vanilla Pythia-2.8B. Given these findings, along with our TL;DR results and the shared reward objective of PPO and Best of 128, we use the Best of 128 baseline as an approximate indicator of PPO-level performance.

In summary, among the methods studied, only DPO delivers consistent gains over the Anthropic HH preferred completions while remaining computationally efficient. It also rivals or exceeds the accuracy of the resource-intensive Best of 128 approach. As shown in Figure~\ref{fig:dialogue-main}, DPO rapidly reaches its optimal performance level during training.

\subsection{Generalization to a new input distribution}

To more thoroughly assess the robustness of PPO and DPO models under shifts in input distribution, we test the policies developed during the Reddit TL;DR summarization experiments on a new domain: news articles from the test set of the CNN/DailyMail corpus \citep{nallapati-etal-2016-abstractive}. We use the optimal sampling temperatures identified for TL;DR (0 and 0.25). The outcomes are displayed in Table~\ref{tab:ood}. To evaluate summary quality, we calculated the GPT-4 win rate in comparison to reference summaries from the dataset, using an adapted GPT-4 (C) prompt from the Reddit TL;DR setup, modifying “forum post” to “news article.” Within this shifted domain, DPO consistently surpasses PPO by a notable margin. These findings suggest that DPO can match or exceed PPO in out-of-distribution generalization performance, even though DPO does not benefit from the extra unlabeled Reddit TL;DR prompts that are available to PPO.

\subsection{Validating GPT-4 judgments with human judgments}
\label{sec:human-judgments}
A human evaluation study was conducted to assess the dependability of GPT-4’s summary judgments using the TL;DR summarization experiment outputs and two types of GPT-4 prompts. The \textbf{GPT-4 (S)} prompt (simple) requests GPT-4 to select the summary that better captures the key information from the post. The \textbf{GPT-4 (C)} prompt (concise) similarly asks for the most concise summary, addressing our observation that with the \textbf{GPT-4 (S)} prompt, GPT-4 tends to favor lengthier, repetitive summaries in contrast to human preferences. The full versions of these prompts can be found in Appendix~\ref{app:prompts}. We include three comparisons in our study: the top-performing method (DPO, temp. 0.25), the lowest-performing method (PPO, temp. 1.0), and an intermediate method (SFT, temp. 0.25), all contrasted with greedily-decoded PPO (its optimal temperature), in order to span a range of summary qualities. The agreement rate between GPT-4 and human raters, with both prompt styles, is comparable to inter-human agreement, implying that GPT-4 can effectively approximate human evaluation. Due to limited availability of human raters, we gathered multiple human assessments only for the DPO and PPO-1 settings. Across all cases, the \textbf{GPT-4 (C)} prompt produces win rates more closely aligned with human judgments, motivating its use for the principal evaluations in Section~\ref{sec:dpo-real-datasets}. Further details about the design of the human evaluation, including the user interface and the identities of volunteers, are available in Appendix~\ref{app:human-study}.

\section{Discussion}
Learning from preferences offers an effective and scalable approach to building robust and aligned language models. In this work, we have presented DPO, a straightforward method for preference-based language model training that does not rely on reinforcement learning. Instead of adapting the preference learning challenge to a classical RL context to utilize standard RL tools, DPO leverages an explicit correspondence between reward functions and language model policies. This framework enables direct alignment with human preferences via a basic cross-entropy objective, circumventing the need for reinforcement learning while preserving generality. Remarkably, DPO achieves comparable or superior results to prevalent RLHF techniques such as those utilizing PPO, often without requiring extensive hyperparameter optimization. As such, DPO substantially lowers the practical challenges associated with training language models from preference data.

\textbf{Limitations \& Future Work.} Our findings open a number of avenues for further exploration. One key question concerns how the generalization capabilities of DPO-trained models compare, particularly out of distribution, to those trained with an explicit reward model. Preliminary experiments indicate that DPO-trained policies may generalize at a level similar to PPO-trained counterparts, yet systematic investigations remain necessary. It is also worth exploring whether techniques such as self-labeling with the DPO policy can effectively leverage unlabeled data. Additionally, understanding the manifestation of reward over-optimization within the DPO framework—and whether the marginal performance drop seen in Figure~\ref{fig:dialogue-main}-right exemplifies this phenomenon—presents another important research direction. Our current assessments focus on models up to 6B parameters; thus, examining how DPO performs when scaled to substantially larger, state-of-the-art architectures represents an intriguing prospect. In terms of evaluation methodology, we observed that GPT-4-derived win rates are influenced by prompt construction, indicating that improved strategies for automated evaluation should be developed. Finally, DPO’s utility is not confined to aligning language models from human preferences—it may also prove valuable for training generative systems in diverse domains beyond text.